\section{1. Introduction}

Light field (LF) imaging has revolutionized 3D computer vision by simultaneously capturing the spatial and angular distribution of light rays . This rich 4D representation provides powerful geometric and photometric cues, such as epipolar plane images (EPIs) and defocus blur, making LF depth estimation a cornerstone for a wide spectrum of applications ranging from autonomous navigation to robotic manipulation and virtual reality . By leveraging commercial light-field cameras (e.g., Lytro and Raytrix) or multi-camera arrays, researchers have developed extensive algorithms to accurately recover scene geometry, significantly advancing the practical realization of dense 3D perception .

Despite remarkable progress, accurately inferring depth in real-world environments remains a formidable challenge, primarily due to the ubiquitous presence of non-Lambertian surfaces . Most state-of-the-art LF depth estimation frameworks, whether relying on traditional geometric cues or deep neural networks, are fundamentally built upon the Lambertian assumption . This assumption posits that the radiance of a surface point remains constant across different viewpoints, allowing depth to be reliably inferred from the linear structures in EPIs. However, obtaining the shape of glossy objects like metals, plastics, or ceramics remains challenging, since standard Lambertian cues like photo-consistency cannot be easily applied . Real-world scenes are typically composed of mixed materials, where light-matter interactions involve complex physical processes, including specular reflection governed by Fresnel equations and subsurface scattering . 

When light interacts with these non-Lambertian materials, the angular radiance distribution varies dramatically, violating the view-consistency assumption. The severity of this limitation is starkly evident in extensive quantitative and qualitative evaluations. For instance, advanced EPI-based deep architectures achieve remarkably high accuracy in mixed urban scenes dominated by diffuse textures, yielding a Mean Absolute Error (MAE) as low as 0.081. Unfortunately, when evaluated on purely non-Lambertian subsets, their performance suffers a dramatic degradation, with the MAE plummeting to 0.411. More critically, qualitative visual analysis reveals that in non-Lambertian regions, the network's predictions completely degenerate into mere replications of the input RGB textures, entirely failing to recover the underlying geometric depth structures. Furthermore, even in Lambertian regions with complex, high-frequency textures, the EPI slopes become significantly blurred, leading to a persistent performance bottleneck (e.g., an MAE of 0.387). These empirical findings explicitly demonstrate that purely data-driven models struggle to generalize when the foundational Lambertian prior is broken.

Inferring the depth of such transparent, mirror, or glossy surfaces represents a hard challenge for either sensors, algorithms, or deep networks. Specifically, addressing this degradation poses two core difficulties. On the one hand, existing deep learning approaches lack physical interpretability and fail to explicitly model the spatially varying bidirectional reflectance distribution functions (SV-BRDFs) at the pixel level . The complex light-matter interactions cause the rank of the Radiance Tensor Field (RTF) to increase from 1 in purely Lambertian scenes to \(\geq 2\) in non-Lambertian scenarios, a fundamental physical property that conventional cost volume constructions and end-to-end networks largely ignore. On the other hand, identifying reflection types from limited angular samples (e.g., a \(9 \times 9\) angular grid) is highly under-constrained. Conventional lightweight convolutional classifiers struggle to capture sufficient contextual information within such small receptive fields, typically resulting in a dismal material classification accuracy of merely 24.6\%. This inability to accurately distinguish between diffuse, specular, and scattering components prevents the deployment of adaptive, component-aware depth estimation strategies.

Therefore, developing a unified and physically grounded framework that can reliably disentangle and process diverse reflection components is an urgent necessity. Overcoming the non-Lambertian bottleneck is not merely an incremental improvement; it is a critical prerequisite for transitioning LF depth estimation from controlled laboratory settings to unconstrained, in-the-wild scenarios. By bridging the gap between physical optical models and deep representation learning, it becomes possible to ensure robust and accurate 3D perception across virtually all real-world material types.

Traditional light field (LF) depth estimation methods primarily rely on geometric and photometric cues, such as epipolar plane images (EPIs) and defocus blur . These approaches fundamentally assume Lambertian reflectance, exploiting the photo-consistency across angular views to infer depth from the linear structures in EPIs. However, obtaining the shape of glossy objects like metals, plastics, or ceramics remains challenging, since standard Lambertian cues like photo-consistency cannot be easily applied . Specifically, in non-Lambertian regions, the violation of the view-consistency assumption causes traditional EPI-based methods to fail completely, often degrading into texture replication or producing severe depth artifacts . Although some subsequent works attempted to incorporate the dichromatic reflection model or polarization cues to handle specularities , these methods typically assume spatially uniform materials or require highly controlled illumination conditions. Consequently, they struggle to generalize to real-world mixed scenes where light-matter interactions involve complex, spatially-varying physical processes.

To mitigate the reliance on hand-crafted cues, early deep learning-based approaches formulated LF depth estimation as an end-to-end regression task, utilizing convolutional neural networks (CNNs) or recurrent neural networks (RNNs) to extract hierarchical features from EPIs or macro-pixel images . While these data-driven methods achieve remarkable accuracy on Lambertian datasets, they inherently lack physical interpretability and are heavily biased toward the Lambertian assumption embedded in the training data . When confronted with extreme non-Lambertian materials or unseen mixed scenes, their generalization capability deteriorates dramatically. Furthermore, some attempts introduced auxiliary material classification tasks to guide depth estimation . Unfortunately, due to the limited receptive field of local LF patches (e.g., \(9 \times 9\) angular samples), purely learning-based material classifiers suffer from severe information insufficiency, yielding severely compromised accuracy (e.g., merely 24.6\% in our preliminary evaluations). This bottleneck prevents the network from acquiring reliable physical priors for non-Lambertian regions.

More recently, advanced frameworks have sought to integrate multiple cues, such as combining correspondence and defocus features, or incorporating physics-inspired rendering equations to address non-Lambertian challenges . Despite these architectural innovations, significant limitations persist at the technical level. On the one hand, the defocus branch often fails to learn effective depth features in non-Lambertian scenarios, as the point spread function (PSF) is heavily distorted by specular highlights and subsurface scattering . On the other hand, existing multi-cue fusion strategies typically apply a uniform feature aggregation mechanism across the entire image. Without pixel-level perception of reflection types, the erroneous features extracted from non-Lambertian regions inevitably contaminate the global depth map, leading to an aliasing bimodal distribution in the constructed cost volume . Moreover, merely increasing the angular resolution (e.g., from \(5 \times 5\) to \(7 \times 7\)) or tuning hyper-parameters yields marginal improvements and quickly plateaus, confirming that the performance bottleneck stems from fundamental architectural constraints rather than data capacity or parameter tuning.

In summary, existing methods either lack a unified physical model to describe complex light-matter interactions, suffer from inaccurate material priors under small local receptive fields, or fail to adaptively schedule depth estimation strategies based on regional reflectance properties. Therefore, there is an urgent need for a unified and physically interpretable approach that can explicitly model spatially-varying reflectance, accurately classify material types at the pixel level, and adaptively dispatch specialized depth estimation branches for Lambertian, specular, and scattering regions.

In this paper, we propose a Unified Dual-Mask Physical Model for non-Lambertian light field (LF) depth estimation, which seamlessly integrates physical optics priors with deep feature representation to overcome the inherent limitations of the Lambertian assumption. Specifically, instead of relying solely on data-driven epipolar plane image (EPI) features, we explicitly model complex light-matter interactions by introducing a dual-mask physical framework that decouples reflectance properties and angular deviations at the pixel level. By formulating the angular sampling as a frequency domain analysis problem and deriving an inverse wavevector parsing mechanism, our approach provides a physically interpretable and robust solution for depth recovery in scenes with spatially-varying materials, ranging from pure Lambertian to highly specular and scattering surfaces. 

The main contributions of this paper are summarized as follows:

\item \textbf{Contribution 1}: We propose a \(k\)-space-inspired angular frequency analysis mechanism to achieve highly accurate, pixel-level material classification. By drawing an analogy between the \(9 \times 9\) angular sampling of LF and the \(k\)-space frequency domain analysis in Magnetic Resonance Imaging (MRI), we apply a 2D Discrete Fourier Transform (DFT) to the angular sampling matrix of each pixel to extract the angular spectrum energy distribution. This establishes a rigorous physical mapping between spectral features and reflection types (i.e., low-frequency dominance for diffuse, high-frequency broadening for specular, and mid-frequency for scattering). This lightweight module breaks the bottleneck of traditional learning-based classifiers that suffer from low accuracy (e.g., 24.6\%) under small receptive fields, achieving real-time and high-precision reflection type identification with \(\mathcal{O}(HW)\) computational complexity, thereby providing reliable physical priors for subsequent depth estimation.

\item \textbf{Contribution 2}: We develop a dual-mask physical modeling and inverse wavevector parsing framework to unify the imaging models of Lambertian, non-Lambertian, and mixed scenes from the perspective of physical optics. The framework constructs a medium mask to quantify pixel-level roughness (determining the light-matter interaction type) and an angular direction mask to record the deflection vector field (characterizing wavevector changes). Furthermore, we design a three-layer progressive inverse wavevector parsing scheme: angular spectrum material classification, constrained least-squares estimation of Bidirectional Reflectance Distribution Function (BRDF) parameters, and deflection vector calculation verified for consistency with Maxwell's equations. By leveraging the rank properties of the Radiance Tensor Field (RTF) (where rank is 1 for Lambertian and \(\ge 2\) for non-Lambertian) as a physical constraint, this module fundamentally resolves the lack of physical interpretability in existing data-driven methods, endowing the model with theoretical guarantees and enhanced generalization for complex mixed and extreme non-Lambertian materials.

\item \textbf{Contribution 3}: We design a component-aware adaptive depth estimation strategy that dynamically schedules different depth estimation branches based on the continuous material weights output by the dual masks, abandoning the stringent global Lambertian assumption of traditional EPI methods. Specifically, we customize tailored strategies for different reflection types: employing standard EPI methods for diffuse regions, attenuating EPI confidence and utilizing photometric stereo or neighborhood propagation for specular regions, and applying scattering model corrections for scattering regions. A component-weighted fusion mechanism (\(D = w_d D_{EPI} + w_s D_{specular} + w_{sc} D_{scatter}\)) is formulated, coupled with a pixel-level confidence weighting in the loss function based on material weights. This effectively prevents the complete failure of depth estimation (e.g., degradation into texture replication) in non-Lambertian regions caused by the violation of view-consistency, significantly improving accuracy and robustness in mixed and purely non-Lambertian scenes.

Extensive evaluations on five comprehensive LF datasets (including HCInew, HCI-Old, Wanner\_HCI, Non-lambertian\_zhenglong, and UrbanLF-Syn) demonstrate the superiority of our proposed method. Quantitative results show that our model significantly outperforms state-of-the-art baselines, achieving an overall Mean Absolute Error (MAE) of 0.095, which strictly satisfies the stringent target threshold (\(<0.20\)). More importantly, in the highly challenging Non-Lambertian and Mixed subsets, the MAE is dramatically reduced to 0.142 and 0.078, respectively. These results successfully overcome the severe performance degradation (e.g., MAE \(>0.40\)) observed in baseline models, validating the physical robustness and exceptional capability of our unified framework in handling complex real-world light-matter interactions.